% =============================================================================
% iMED Challenge 2026 -- Structured Method Report
% =============================================================================
%
% This report will be used as source material for the joint iMED Challenge paper.
% Please keep the section/subsection structure unchanged so that reports from
% different teams can be combined automatically using PyLaTeX.
%
% Instructions:
% - Replace every \TODO{...}. Write "N/A" where a field does not apply.
% - Complete only the task(s) your team submitted to.
% - Describe the FINAL submitted method.
% - Do not add sections (I will write Introduction, Dataset Desc. etc.).
% - Clearly disclose public external data and pretrained models.
% - Optional figures should be supplied alongside this .tex file.
% =============================================================================
% REFERENCES
% =============================================================================
% Include the complete BibTeX entry for every reference cited in this report.
% Do not submit a separate .bib file.

\begin{filecontents*}[overwrite]{references.bib}

@article{kerbl2023,
  author  = {Kerbl, Bernhard and Kopanas, Georgios and Leimk{\"u}hler, Thomas and Drettakis, George},
  title   = {3D Gaussian Splatting for Real-Time Radiance Field Rendering},
  journal = {ACM Transactions on Graphics},
  year    = {2023},
  volume  = {42},
  number  = {4}
}

@inproceedings{huang2024endo4dgs,
  author    = {Huang, Yiming and Cui, Beilei and Bai, Long and Guo, Ziqi and Xu, Mengya and Islam, Mobarakol and Ren, Hongliang},
  title     = {Endo-4DGS: Endoscopic Monocular Scene Reconstruction with 4D Gaussian Splatting},
  booktitle = {Medical Image Computing and Computer Assisted Intervention (MICCAI)},
  year      = {2024}
}

@article{zhao2023aliked,
  author  = {Zhao, Xiaoming and Wu, Xingming and Chen, Weihai and Chen, Peter C. Y. and Xu, Qingsong and Li, Zhengguo},
  title   = {{ALIKED}: A Lighter Keypoint and Descriptor Extraction Network via Deformable Transformation},
  journal = {IEEE Transactions on Instrumentation and Measurement},
  year    = {2023},
  volume  = {72}
}

@inproceedings{lindenberger2023lightglue,
  author    = {Lindenberger, Philipp and Sarlin, Paul-Edouard and Pollefeys, Marc},
  title     = {{LightGlue}: Local Feature Matching at Light Speed},
  booktitle = {IEEE/CVF International Conference on Computer Vision (ICCV)},
  year      = {2023}
}

@inproceedings{sun2021loftr,
  author    = {Sun, Jiaming and Shen, Zehong and Wang, Yuang and Bao, Hujun and Zhou, Xiaowei},
  title     = {{LoFTR}: Detector-Free Local Feature Matching with Transformers},
  booktitle = {IEEE/CVF Conference on Computer Vision and Pattern Recognition (CVPR)},
  year      = {2021}
}

@article{umeyama1991,
  author  = {Umeyama, Shinji},
  title   = {Least-Squares Estimation of Transformation Parameters Between Two Point Patterns},
  journal = {IEEE Transactions on Pattern Analysis and Machine Intelligence},
  year    = {1991},
  volume  = {13},
  number  = {4},
  pages   = {376--380}
}

@article{telea2004,
  author  = {Telea, Alexandru},
  title   = {An Image Inpainting Technique Based on the Fast Marching Method},
  journal = {Journal of Graphics Tools},
  year    = {2004},
  volume  = {9},
  number  = {1},
  pages   = {23--34}
}

@inproceedings{gortler1996,
  author    = {Gortler, Steven J. and Grzeszczuk, Radek and Szeliski, Richard and Cohen, Michael F.},
  title     = {The Lumigraph},
  booktitle = {SIGGRAPH},
  year      = {1996},
  pages     = {43--54}
}

@inproceedings{ronneberger2015unet,
  author    = {Ronneberger, Olaf and Fischer, Philipp and Brox, Thomas},
  title     = {{U-Net}: Convolutional Networks for Biomedical Image Segmentation},
  booktitle = {Medical Image Computing and Computer-Assisted Intervention (MICCAI)},
  year      = {2015}
}

@inproceedings{liu2022convnext,
  author    = {Liu, Zhuang and Mao, Hanzi and Wu, Chao-Yuan and Feichtenhofer, Christoph and Darrell, Trevor and Xie, Saining},
  title     = {A {ConvNet} for the 2020s},
  booktitle = {IEEE/CVF Conference on Computer Vision and Pattern Recognition (CVPR)},
  year      = {2022}
}

@misc{simeoni2025dinov3,
  author       = {Sim{\'e}oni, Oriane and others},
  title        = {{DINOv3}},
  year         = {2025},
  howpublished = {arXiv preprint},
  note         = {Self-supervised visual features; the ConvNeXt-Base encoder distilled on LVD-1689M is used here via \texttt{timm}}
}

@article{twinanda2017endonet,
  author  = {Twinanda, Andru P. and Shehata, Sherif and Mutter, Didier and Marescaux, Jacques and de Mathelin, Michel and Padoy, Nicolas},
  title   = {{EndoNet}: A Deep Architecture for Recognition Tasks on Laparoscopic Videos},
  journal = {IEEE Transactions on Medical Imaging},
  year    = {2017},
  volume  = {36},
  number  = {1},
  pages   = {86--97},
  note    = {Introduces the public Cholec80 dataset}
}

@misc{iakubovskii2019smp,
  author       = {Iakubovskii, Pavel},
  title        = {Segmentation Models Pytorch},
  year         = {2019},
  howpublished = {\url{https://github.com/qubvel/segmentation_models.pytorch}}
}

@inproceedings{wang2025vggt,
  author    = {Wang, Jianyuan and Chen, Minghao and Karaev, Nikita and Vedaldi, Andrea and Rupprecht, Christian and Novotny, David},
  title     = {{VGGT}: Visual Geometry Grounded Transformer},
  booktitle = {IEEE/CVF Conference on Computer Vision and Pattern Recognition (CVPR)},
  year      = {2025}
}

% NOTE: please verify this entry against the official release before publication.
@misc{depthanything3,
  author       = {{Depth Anything Team}},
  title        = {Depth Anything 3: Recovering the Visual Space from Any Views},
  year         = {2025},
  howpublished = {\url{https://depth-anything-3.github.io/}}
}

\end{filecontents*}

\documentclass[10pt]{article}

\usepackage[a4paper,margin=2cm]{geometry}
\usepackage{booktabs}
\usepackage{tabularx}
\usepackage{array}
\usepackage{graphicx}
\usepackage{amsmath}
\usepackage[hidelinks]{hyperref}
\usepackage{xcolor}

\setlength{\parindent}{0pt}
\setlength{\parskip}{4pt}
\renewcommand{\arraystretch}{1.12}

\newcommand{\TODO}[1]{\textbf{\textcolor{red!70!black}{[TODO: #1]}}}
\newcolumntype{Y}{>{\raggedright\arraybackslash}X}
\newcolumntype{L}[1]{>{\raggedright\arraybackslash}p{#1}}

\begin{document}

\begin{center}
    {\Large\bfseries iMED Challenge 2026}\\
    {\large Structured Method Report}
\end{center}

% =============================================================================
\section*{Submission Information}
% =============================================================================

\begin{tabularx}{\textwidth}{L{0.27\textwidth}Y}
\toprule
Team name & Jmees \\
Method name / acronym & \textbf{Tri3D-GridBA} (PE) and \textbf{SDW} --- Static-rig Depth Warp (NVS) \\
Task(s) submitted & Both (iMED PE and iMED NVS) \\
Primary contact & Shunsuke Kikuchi, \texttt{kikushun2111@gmail.com} \\
Challenge-paper contributors &
\TODO{List up to 3 proposed contributors in priority order (highest priority first). The final number of authors included per team may be reduced depending on the number of eligible reports received.} \\
Affiliation(s) &
\TODO{Full institution name(s); indicate the affiliation(s) corresponding to each proposed contributor.} \\
Public code (optional) & \TODO{Insert the public repository URL} \\
\bottomrule
\end{tabularx}

Final submitted images: \texttt{docker.synapse.org/syn74277461/jmees26-pe:v1} (PE) and
\texttt{docker.synapse.org/syn74277461/jmees26-nvs:v12} (NVS).

% =============================================================================
\section{Methodology}
% =============================================================================

% Describe how the FINAL submitted method works.
% Focus on the conceptual pipeline and the main technical components.
% Do not write implementation settings such as hardware, learning rate,
% batch size, exact training duration, or preprocessing parameters; these belong
% in Section 3.
%
% If you submitted to both tasks, complete both subsections separately.

\subsection{iMED Pose Estimation (PE)}

Our final PE method, \textbf{Tri3D-GridBA}, is a purely geometric (non-learned) pipeline. It keeps the
front end of the official baseline --- ALIKED keypoints~\cite{zhao2023aliked} matched with
LightGlue~\cite{lindenberger2023lightglue} --- and replaces the entire back end.

\paragraph{Problem reading.} The quantity to predict is the \emph{same-time cross-camera} relative pose
$T_{e2\leftarrow e1}(t)$ between the two endoscopes, expressed frame-to-initial
($T(0)^{-1}T(t)$, frame $0$ = identity), and it is scored by Sim(3)-aligned ATE. Two consequences drive
the design: (i) ATE depends only on the \emph{translation} of each pose, and (ii) Sim(3) alignment absorbs
one global scale, so an absolute metric scale is unnecessary but the \emph{relative} scale between frames
must be consistent. The official baseline recovers the cross-camera pose from an essential matrix and
therefore normalises the translation to unit norm, discarding exactly this frame-to-frame scale
information; recovering it is the single largest source of our gain.

\paragraph{Visual information used.} Both stereo pairs are used at every time step. Per frame we compute
three LightGlue matches from a shared ALIKED feature cache: the endoscope-1 stereo pair
$e1_L\!\leftrightarrow\!e1_R$, the endoscope-2 stereo pair $e2_L\!\leftrightarrow\!e2_R$, and the
cross-camera pair $e1_L\!\leftrightarrow\!e2_L$. Because features are extracted once per image and matched
through shared keypoint indices, a keypoint triangulated in the $e1$ stereo pair can be identified with a
keypoint triangulated in the $e2$ stereo pair whenever the cross match links them.

\paragraph{Fixed stereo extrinsics.} Each endoscope is a rigid stereo rig, so its left--right extrinsic is
constant over a sequence. We estimate it once per scope from per-frame essential matrices and aggregate
over \emph{all} frames (rotation averaging, mean unit translation). This removes the per-frame rotation
noise of a single-frame estimate and, more importantly, fixes one constant triangulation scale per scope
for the whole sequence.

\paragraph{Per-frame pose by robust 3D--3D alignment.} With the fixed extrinsics we triangulate two point
clouds per frame, $\mathbf{X}^{(1)}$ in the $e1_L$ frame and $\mathbf{X}^{(2)}$ in the $e2_L$ frame, and
pair them through the cross match. The relative pose is the Sim(3) transform that aligns them, solved in
closed form by a weighted Umeyama estimator~\cite{umeyama1991},
\begin{equation}
  (s_t,\mathbf{R}_t,\mathbf{c}_t)\;=\;\arg\min \sum_i w_i\,\bigl\| s\,\mathbf{R}\,\mathbf{x}^{(1)}_i + \mathbf{c} - \mathbf{x}^{(2)}_i \bigr\|^2 ,
\end{equation}
wrapped in Tukey-biweight IRLS. The robustifier is essential and its interpretation is specific to this
task: for a \emph{same-time} rigid cross-camera pose, any correct correspondence already determines the
transform exactly, so every point that disagrees with the dominant model is a mismatch, a
non-rigidly-moving point (tissue deformation, instrument), or a point corrupted by inter-camera
synchronisation jitter. All three are handled as outliers. Weighting by an independent cue (image
velocity from dense optical flow) was implemented and tested but \emph{hurt}; the fit residual itself is
the better signal, so the deployed system uses residual-driven weights only.

\paragraph{Grid bundle adjustment.} The per-frame estimates are independent and therefore jittery on
sequences with fast motion. We add temporal constraints without giving up the per-frame data. If
$M_{e1}(t)$ and $M_{e2}(t)$ are the ego-motions of each scope between $t$ and $t{+}1$, the cross-camera
poses obey exactly
\begin{equation}
  T(t{+}1) \;=\; M_{e2}(t)\, T(t)\, M_{e1}(t)^{-1}.
\end{equation}
We estimate $M_{e1},M_{e2}$ by the same robust 3D--3D alignment applied \emph{temporally} within each
scope, reusing the per-frame stereo clouds so only the temporal $L\!\leftrightarrow\!L$ match is new.
Holding the rotations fixed at their per-frame estimates, the relation becomes \emph{linear in the
translations}, $\mathbf{c}(t{+}1) = \mathbf{B}(t)\,\mathbf{c}(t) + \mathbf{d}(t)$ with
$\mathbf{B} = \mathbf{R}_{M_{e2}}$, so absolute (same-time) terms, relative (ego-motion) terms and a
second-difference smoothness prior can be fused in one linear least-squares problem:
\begin{equation}
  \min_{\{\mathbf{c}(t)\}} \; \sum_t w_t \bigl\|\mathbf{c}(t)-\mathbf{c}^{\mathrm{meas}}(t)\bigr\|^2
  \;+\; \lambda_{\mathrm{rel}} \sum_{(t,u)\in\mathcal{E}} \bigl\|\mathbf{c}(u)-\mathbf{B}\mathbf{c}(t)-\mathbf{d}\bigr\|^2
  \;+\; \lambda \sum_t \bigl\|\mathbf{c}(t{+}1)-2\mathbf{c}(t)+\mathbf{c}(t{-}1)\bigr\|^2 .
\end{equation}
Only translations are optimised, since ATE does not see rotation. The absolute weight $w_t$ is a
per-frame confidence built from the inlier count and the geometric spread of the 3D correspondences,
divided by the fit residual: an ill-conditioned, clustered point set yields low confidence and is pulled
toward its better-conditioned neighbours, while a well-conditioned frame is left essentially untouched.
The ego-motion term carries a deliberately weak weight because scene deformation makes the rigid
ego-motion assumption only approximate.

\paragraph{Scale and failed registrations.} No absolute scale is estimated. One constant stereo baseline
per scope gives a single, self-consistent scale for the whole sequence, which is all the Sim(3)-aligned
metric requires. Frames with too few matches or too few valid 3D--3D pairs are marked unregistered; the
grid BA still produces a translation for them through the smoothness and ego-motion terms, so the output
is dense (100\% of frames written). If fewer than three frames register at all, the system falls back to
the plain per-frame path. Finally, after the BA the whole track is re-anchored so that frame $0$ is
exactly the identity, as the output contract requires --- a pure translation shift that the Sim(3)-aligned
metric absorbs and that leaves the trajectory shape unchanged.

\paragraph{Intrinsics guard.} The released PE intrinsics were, at one point, distributed at the native
calibration resolution rather than the frame resolution. The submitted container detects this case from
the principal point ($c_x > 0.65\,W$) and rescales $K$ by $W/1280$; an already-corrected $K$ passes
through unchanged. This makes the container correct under either version of the data, and it is the only
difference between our final image and the two other PE images we submitted.

\paragraph{Summary of the modifications to the base method.} Relative to the official baseline we keep
ALIKED+LightGlue and change: (i) 2D--2D essential matrix $\rightarrow$ two-sided 3D--3D Sim(3) alignment
with a consistent stereo scale; (ii) unit translation $\rightarrow$ metric-consistent translation;
(iii) per-frame independent poses $\rightarrow$ frame-to-initial parameterisation plus grid BA over
ego-motion edges; (iv) single-frame stereo extrinsics $\rightarrow$ all-frame aggregation;
(v) least squares $\rightarrow$ Tukey IRLS.

\subsection{iMED Novel View Synthesis (NVS)}

Our final NVS method, \textbf{SDW} (Static-rig Depth Warp), performs \emph{no per-sequence optimisation at
all}: each target frame is rendered by a single feed-forward geometric warp of the corresponding source
frame, followed by hole filling and a fixed photometric transform.

\paragraph{Problem reading.} \texttt{pose.txt} contains exactly two entries per sequence --- endoscope 2
at the identity and endoscope 1 at a fixed relative pose. The two cameras therefore do not move; only the
scene deforms. The task is not novel-view synthesis along a trajectory but a per-frame re-projection
between two \emph{static} viewpoints of a deforming scene. Consequently the deformation does not need to
be modelled at all: it is fully observed in the per-frame source image and its per-frame depth map. This
observation is what allows a training-free warp to match a 4D Gaussian Splatting
baseline~\cite{huang2024endo4dgs,kerbl2023} that optimises for hours per sequence.

\paragraph{Representation.} The scene representation is the per-frame organiser-supplied depth map of
\texttt{endoscope2/L}, i.e.\ a per-frame point cloud, rendered by forward splatting with a $z$-buffer.
There is no persistent scene model, no radiance field and no learned appearance.

\paragraph{Pipeline.} For each frame:
\begin{enumerate}\itemsep2pt
  \item \textbf{Instrument exclusion.} A surgical-instrument segmenter (ConvNeXt-Base
        U-Net~\cite{liu2022convnext,ronneberger2015unet}) is run on the source image and the depth of the
        detected instrument pixels is zeroed so they are never splatted. The reason is a property of the
        official metric: instrument pixels in \emph{endoscope 1} are excluded from the scored region, so a
        re-projected instrument can only land where the ground truth shows \emph{tissue}, i.e.\ it paints
        metal over organ inside the scored area. The measured error profile confirms this: reconstruction
        error rises monotonically toward the instrument mask (18.7\,dB within 2\,px of it versus 22.1\,dB
        beyond 64\,px). We use a segmentation model rather than the supplied \texttt{toolL} masks because
        the supplied masks are not reliable (on two sequences they paint a persistent blob where no
        instrument exists, including frames in which the instrument has left the field of view), and
        because the model needs nothing but the source image, which the hidden test is guaranteed to
        provide.
  \item \textbf{Geometric warp.} Pixels are back-projected with the source intrinsics,
        $\mathbf{X}_2 = D(u,v)\,K_2^{-1}[u,v,1]^{\!\top}$, transformed by the constant relative pose
        $T_{\mathrm{rel}} = T_{e1}^{-1}T_{e2}$, and projected with the target intrinsics $K_1$.
  \item \textbf{Visibility.} Splatting is resolved with a $z$-buffer implemented as a single
        scatter-reduce: the quantised depth and the source pixel index are packed into one 64-bit integer
        so that a minimum-reduction over the target pixel selects the nearest surface. This is exact
        (no blending) and runs entirely on the GPU. Soft/bilinear splatting was tested and is clearly
        worse, because averaging several sources per target pixel injects blur.
  \item \textbf{Hole filling.} Disocclusions --- regions visible to endoscope 1 but occluded in
        endoscope 2 --- are not a rounding detail: they cover $19.6\%$ of the officially scored pixels
        (15--27\% per sequence), and replacing them with ground truth is worth $+0.99$\,dB. They are
        filled with a GPU push--pull pyramid~\cite{gortler1996}. A structure-aware alternative
        (Fast-Marching inpainting~\cite{telea2004}) raises PSNR but costs $0.044$ SSIM and lost on the
        organiser leaderboard, so push--pull is the deployed filler (Section~2.2).
  \item \textbf{Photometric transform.} The two endoscopes differ systematically in white balance and
        response. A fixed per-channel affine tone curve $\hat{y}_c = \mathrm{clip}(g_c x_c + b_c)$ is
        applied to the warp output. Its six constants are fitted once, offline, by least squares over the
        scored region of the 20 public sequences and then frozen; no ground truth is read at inference.
\end{enumerate}

\paragraph{Camera poses.} Poses are used exactly as supplied and are never refined. We measured that a
per-sequence geometric correction would be worth up to $\sim\!2$\,dB, but it does not transfer across
sessions (Section~2.2), so nothing is fitted per scene.

\paragraph{Summary of the modifications to the base method.} Relative to the Endo-4DGS
baseline~\cite{huang2024endo4dgs} the method is a replacement rather than a modification: no 4D Gaussian
representation, no deformation field, no per-sequence training. What we add on top of a plain
depth-reprojection warp is (i) exact GPU $z$-buffered splatting, (ii) instrument exclusion driven by the
scoring mask definition, (iii) push--pull disocclusion filling, and (iv) a frozen low-order photometric
transform.

\section{Experiments Performed}

\subsection{iMED Pose Estimation (PE)}

\paragraph{Development split.} We used the 19 released test sequences that ship with ground-truth
\texttt{pose.txt} as our development set, and grouped sequences by \texttt{session} for
leave-one-session-out (LOSO) evaluation whenever any quantity was fitted, because sequences from the same
session share a scene and leak. The organiser validation leaderboard (6 sequences) was used as a second,
independent readout.

\paragraph{Metric.} Sim(3)-aligned mean ATE in mm, computed with the organiser evaluation script. We
additionally tracked the RMSE variant, because the leaderboard column is
\texttt{mean\_ate\_rmse\_mm}; ablation conclusions were checked under both.

\paragraph{Main ablation.} All numbers are mean ATE over the 19 development sequences.

\begin{center}
\begin{tabular}{lcc}
\toprule
Configuration & mean ATE [mm] & vs.\ baseline \\
\midrule
Official baseline (ALIKED+LightGlue, essential matrix, unit $t$) & 2.540 & --- \\
Tri3D: two-sided 3D--3D Umeyama, consistent stereo scale        & 1.034 & $-59\%$ \\
\quad + Tukey IRLS robust rigid fit                              & 0.998 & $-61\%$ \\
\quad + grid BA ($K=1$, $\lambda_{\mathrm{rel}}=0.3$) \textbf{[submitted]} & \textbf{0.953} & $\mathbf{-62.5\%}$ \\
\bottomrule
\end{tabular}
\end{center}

An earlier ablation, run before the intrinsics were corrected, isolated the two structural ingredients
separately and is reported here because the two are orthogonal: Tri3D alone $1.423$; $+$ all-frame stereo
extrinsic aggregation $1.346$; $+$ temporal BA ($\lambda=2$) $1.292$; both $1.218$.

\paragraph{What did not work.} We tested learned alternatives at every stage and none beat the classical
pipeline.
\begin{itemize}\itemsep1pt
  \item \emph{Learned multi-view pose.} VGGT~\cite{wang2025vggt} applied to the same-time cross-camera
        pair: $1.79$\,mm (better than the baseline, worse than Tri3D). Depth-Anything-3~\cite{depthanything3}
        in the same recipe: $1.39$\,mm versus $1.27$\,mm for Tri3D on the same 6-sequence subset ---
        the best learned variant, still short.
  \item \emph{Learned depth instead of triangulation.} Substituting DA3 metric depth for stereo
        triangulation gave $1.084$ versus $1.014$ for Tri3D (6-sequence subset, both with BA). Treating
        DA3 output as relative depth and affinely fitting it to stereo was much worse ($1.998$ vs
        $1.060$). The conclusion is conceptual: Tri3D already \emph{is} a correctly-scaled relative depth
        estimator, and the Sim(3) metric absorbs the remaining global freedom, so a learned metric depth
        adds nothing it can beat geometry with.
  \item \emph{Learned confidence for the BA.} Per-frame features do correlate with error
        (e.g.\ cross-match count, $r=-0.50$), but a learned weighting was worse than the analytic one
        under LOSO ($1.092$ vs $1.067$) and on test ($1.266$ vs $1.218$).
  \item \emph{Other matchers.} On a 6-sequence subset, ALIKED $1.014$ beat DoG+HardNet $1.191$,
        SuperPoint $1.223$, DISK $1.245$ and dense LoFTR~\cite{sun2021loftr} $1.176$. Match
        \emph{precision}, not match count, is what this task rewards.
  \item \emph{Ensembling.} A Tri3D$\times$VGGT oracle was worth only $-1.9\%$, so an ensemble was not
        pursued in the final image. A three-way ensemble with a windowed VGGT variant was built and
        measured: on the hard subset, an exhaustive weight search assigned VGGT a weight of \emph{exactly
        zero} in all three variants we tried. Its errors are correlated with Tri3D's because they share
        the same observability limit, so it contributes no independent information where it would matter.
  \item \emph{Temporal median smoothing} reduced the mean by $5\%$ but degraded 6 of 19 sequences,
        including the best-behaved zoom-in trajectories; it is off. \emph{fp16 matching} degraded ATE from
        $0.953$ to $1.05$ and is off.
\end{itemize}

\paragraph{Diagnosis of the residual error.} The worst sequences are the fast circular trajectories of two
particular sessions. We eliminated the obvious causes one by one: correspondence count (the worst
sequences have \emph{more} matches, $\sim$630), scale noise (CV $0.02$), registration failures (none), and
depth quality (disparity 23.7\,px, Umeyama residual 0.012). What remains is that geometrically well-fitted
frames still show roughly $2\times$ translation error under fast lateral motion --- an ill-conditioned
point distribution combined with inter-camera synchronisation, i.e.\ an observability limit of the data
rather than a defect of the estimator. This is consistent with four independent methods failing on the
same sequences.

\paragraph{Final configuration.} Tri3D $+$ all-frame stereo aggregation $+$ IRLS $+$ grid BA
($K=1$, $\lambda=2.0$, $\lambda_{\mathrm{rel}}=0.3$), with the intrinsics guard enabled. Among the three
PE images we submitted (which differ \emph{only} in intrinsics handling and in whether a DA3 blend is
applied), we selected the one that auto-corrects the intrinsics, because it is the only image that is
correct under both versions of the released calibration; its leaderboard rank and spacing also matched the
prediction from our local hard-subset measurements. We note explicitly that the leaderboard differences
between our three PE images ($2.296$/$2.323$/$2.345$) are inside the resolution of a 6-sequence set: a
bootstrap over paired per-sequence differences (20{,}000 resamples of 6 from our 19) gives $\pm 0.10$\,mm,
so those gaps were not treated as evidence.

\paragraph{Guarding against leaderboard overfitting.} We used neither readout as a search objective. All
hyperparameters ($\lambda$, $\lambda_{\mathrm{rel}}$, the BA window $K$, the robustifier, the aggregation
count) were fixed on the 19-sequence development set before our first submission and never revised on
leaderboard feedback; our three submitted images differ only in intrinsics handling and an auxiliary blend,
neither of which is a tuned scalar. At $\pm 0.10$\,mm resolution, tuning against six sequences would be
fitting noise.

\subsection{iMED Novel View Synthesis (NVS)}

\paragraph{Development split.} All 20 released \texttt{iMED\_NVS} sequences (3{,}953 frames) with
endoscope-1 ground truth. Anything fitted (the photometric constants) was evaluated leave-one-session-out
over the 7 sessions; renderer-level A/B comparisons were run over all 20 sequences. The organiser
validation leaderboard was used as an independent readout. Two measurement systems appear below and are
\emph{not} directly comparable: an offline LOSO protocol on a frame subsample, and whole-set runs of the
actual containers; comparisons are always within one system.

\paragraph{Metric.} The official \texttt{metrics.py} (PSNR primary; SSIM and LPIPS secondary), evaluated
inside the official region (endoscope2$\rightarrow$endoscope1 overlap $\times$ instrument-free). We
verified that the metric implementation inside the official baseline container is byte-identical to the
released one, and re-derived the effective scoring region from it directly, because self-evaluation over
the full frame under-reports by several dB.

\paragraph{Warp versus the 4DGS baseline.} On the one sequence for which we ran the official baseline to
convergence (\texttt{session\_004\_scene\_2\_tool\_1}, native resolution):

\begin{center}
\begin{tabular}{lcccl}
\toprule
Method & PSNR & SSIM & LPIPS & Cost \\
\midrule
Official Endo-4DGS (coarse 2000 + fine 6000 it.) & 21.14 & 0.699 & 0.0190 & $\sim$2\,h optimisation/seq \\
Ours, depth-reprojection warp                    & \textbf{21.45} & \textbf{0.721} & 0.0198 & inference only \\
\bottomrule
\end{tabular}
\end{center}

\paragraph{Photometric transform family.} Fitting the transform is the single largest gain, and its
\emph{capacity} matters more than its form (20 sequences, LOSO by session, official mask):

\begin{center}
\begin{tabular}{lcccccc}
\toprule
Transform & identity & affine (6) & poly2 (9) & poly3 & $3{\times}3$ matrix & LUT-32 \\
\midrule
PSNR & 19.934 & 22.022 & 22.149 & 22.021 & 20.927 & 21.597 \\
\bottomrule
\end{tabular}
\end{center}

Spatially-varying variants ($21.776$) and a root-polynomial variant also lost. The pattern is consistent:
only a low-order model that corresponds to actual camera physics transfers across sessions; every extra
degree of freedom turns into per-scene overfitting.

\paragraph{Renderer factorial.} After a submission that bundled two changes and lost, we made every knob
an environment variable and ran a $2\times 2$ factorial from a \emph{single} image over all 20 sequences /
3{,}953 frames:

\begin{center}
\begin{tabular}{llcc}
\toprule
Instrument excluded & Hole filler & PSNR & SSIM (gray) \\
\midrule
no  & push--pull & 22.3625 & 0.6892 \\
\textbf{yes} & \textbf{push--pull} & \textbf{22.4880} & \textbf{0.6939} \\
no  & Fast-Marching inpainting & 22.4473 & 0.6449 \\
yes & Fast-Marching inpainting & 22.5602 & 0.6499 \\
\bottomrule
\end{tabular}
\end{center}

The main effects are additive and consistent at both levels of the other factor: instrument exclusion is
$+0.126$ PSNR / $+0.005$ SSIM (a clean win on both metrics), whereas inpainting is $+0.085$ PSNR but
$-0.044$ SSIM. The bundled submission had lost because of the inpainting, not the instrument mask --- which
we had nearly discarded along with it. This is why the deployed filler is push--pull despite its lower
PSNR. A push--pull pyramid-depth sweep ($3$: 21.43, $4$: 21.87, $6$: 22.49, $8$: 22.59, $10$: 22.59,
saturating at $12$) selected a deeper pyramid on our development set, but that change \emph{lost} on the
organiser leaderboard, so the submitted image keeps the default depth.

\paragraph{Where the remaining error is.} We decomposed the residual inside the official mask. It is
dominated by per-\emph{scene} calibration mismatch: a sequence-constant global image shift of 8--35\,px
(intensity-based global alignment is worth $+1.89$\,dB; fitting a Sim(3) pose correction is worth
$+1.99$\,dB), and a photometric offset whose oracle gain of $+1.44$\,dB is $\sim$90\% scene-level and only
$+0.14$\,dB frame-level. Combining both oracles reaches $25.879$ versus $21.917$ deployed. Crucially,
\emph{both corrections fail to transfer under LOSO} (a burned-in pose correction costs $-0.05$ to
$-0.40$\,dB on unseen sessions; a photometric transform predicted from source-image statistics scores
$20.393$, winning on 3 of 20). Estimating them requires observing endoscope 1, which is not permitted.
We therefore conclude that our renderer is close to the ceiling of what is observable from endoscope 2
alone, and we deliberately stopped adding per-scene machinery.

\paragraph{Other rejected ideas.} Multi-source temporal warping (filling holes from neighbouring frames
$t\pm k$) was worse than single-frame inpainting on all four sequences tested --- in a deforming scene a
neighbouring frame supplies misaligned texture. Bilinear/softmax splatting cost $2.29$\,dB. Restricting
the inpainting neighbourhood cost $0.5$\,dB. A small U-Net refiner ($+1.59$\,dB) was beaten by the
six-constant affine transform ($+1.83$\,dB in the same protocol).

\paragraph{Final configuration and one deliberate exception.} The submitted image is: GPU warp $+$
instrument exclusion $+$ push--pull filling $+$ \emph{affine} photometric transform. The last choice is
made \emph{against} our development CV, which prefers the 9-constant poly2 by $0.23$\,dB, and \emph{for}
the leaderboard. The reason is that two independent leaderboard pairs of containers differing in nothing
but this transform gave the same result twice ($+0.058$ and $+0.057$ in favour of affine), while every
parameter we had selected by maximising development CV had inverted on the held-out set (the photometric
order, the hole filler, the pyramid depth --- three of four), and the one change that did transfer was the
one derived from the geometry of the scoring mask rather than from a CV sweep. This is the same
``only low-order models transfer'' result as the transform-family sweep above.

\paragraph{Guarding against leaderboard overfitting.} The two readouts disagree systematically --- three of
the four settings we chose by maximising CV inverted on the held-out set --- so we treated neither as a
search objective. Each submission therefore changed exactly one factor, and the only leaderboard-driven
decision was a binary choice between two orders of the same photometric model, fixed by a decision rule
written down before the submission was scored and accepted only because the effect had already replicated
at the same magnitude in an independent pair of containers. No continuous parameter was tuned on
leaderboard feedback: the one sweep we ran (pyramid depth) was CV-driven, lost, and was reverted to its
default rather than re-tuned, after which we stopped parameter search.

% =============================================================================
\section{Implementation Details}
% =============================================================================

% Report concrete implementation settings for the FINAL submitted method.
% Do not repeat the method description from Section 1 or the experimental
% comparisons from Section 2.
%
% Clearly disclose all pretrained models/checkpoints and public external data.

\subsection{iMED Pose Estimation (PE)}

\paragraph{Hardware.} Development on a single NVIDIA Quadro RTX 8000 (48\,GB); one GPU is used at
inference. The evaluation host is a single RTX 4090 with a 10-minute budget for all sequences combined.

\paragraph{Training.} \textbf{None.} The submitted PE method contains no trained component and no fitted
constants; it is purely geometric. Consequently there is no optimiser, learning rate, batch size, training
schedule, loss weighting or data augmentation to report.

\paragraph{Pretrained models and external data.} ALIKED (\texttt{aliked-n16}) and LightGlue
(\texttt{aliked} weights), both from the public \texttt{cvg/LightGlue} release and used off the shelf with
no fine-tuning. The weights are baked into the image at build time so the container runs with
\texttt{-{}-network=none}. \textbf{No external datasets were used}, and no iMED data was used to fit
anything.

\paragraph{Input and preprocessing.} Frames are consumed at their native resolution (600$\times$480 in the
released PE data) with no resizing, cropping or augmentation. Intrinsics are read from \texttt{K.txt};
if the principal point indicates an unscaled native-resolution calibration ($c_x > 0.65\,W$), $K$ is
rescaled by $W/1280$, otherwise it is used verbatim.

\paragraph{Feature extraction and matching.} ALIKED with \texttt{max\_num\_keypoints}\,=\,2048; LightGlue
with default thresholds; fp32 throughout (fp16 was measured to degrade ATE from 0.953 to 1.05\,mm).
Features are extracted once per image and cached; the cache is shared by the stereo-extrinsic aggregation,
the same-time cross-camera matches and the temporal ego-motion edges.

\paragraph{Geometric parameters.}
\begin{itemize}\itemsep1pt
  \item Minimum matches to attempt a frame: 8; minimum 3D--3D correspondences to accept a pose: 12.
  \item Stereo extrinsic: estimated per frame by \texttt{findEssentialMat} (RANSAC, prob 0.999,
        threshold $10^{-3}$ on normalised coordinates) and aggregated over \emph{all} frames
        (\texttt{-{}-nagg 9999}): rotation averaging plus mean unit translation.
  \item Robust rigid fit: Tukey biweight IRLS, tuning constant $4.685\,\sigma$ with
        $\sigma$ = median residual, 5 reweighting iterations. RANSAC on the rigid model is available but
        disabled. Velocity-based prior weights are disabled ($\alpha = 0$).
  \item Grid BA: temporal window $K=1$; smoothness weight $\lambda = 2.0$; relative (ego-motion) weight
        $\lambda_{\mathrm{rel}} = 0.3$, with each edge scaled by $\min(n_1,n_2)/50$ where $n_i$ are the
        temporal inlier counts; absolute weight $w_t = n_t \cdot \mathrm{spread}_t / (1 + 5\,\mathrm{resid}_t)$
        normalised by its median over valid frames, where $\mathrm{spread}$ is the median distance of the
        3D correspondences from their median. The normal matrix is regularised with $10^{-6}\mathbf{I}$
        to fix the gauge across gaps. Ego-motion edges use IRLS with the same settings.
  \item Post-processing: frame-to-initial, then an exact re-anchoring of frame 0 to
        $(\mathbf{t}=\mathbf{0},\,q=\mathrm{identity})$.
\end{itemize}

\paragraph{Runtime.} $\approx$0.43\,s per frame; 11 minutes for all 19 development sequences on the
RTX 8000, i.e.\ comfortably inside the budget on the faster evaluation GPU. The feature cache and the
reuse of per-frame stereo clouds by the ego-motion edges cut this from 18.5 minutes with bit-identical
output.

\paragraph{Container.} \texttt{nvidia/cuda:11.8.0-cudnn8-runtime-ubuntu22.04}; Python 3.10;
\texttt{torch 2.1.2+cu118}, \texttt{torchvision 0.16.2+cu118}, \texttt{numpy 1.26.4},
\texttt{scipy 1.11.4}, \texttt{opencv-python-headless 4.9.0.80}, LightGlue from source. Image size
9.37\,GB. Contract: \texttt{/input} mounted read-only, one
\texttt{/output/<seq>/pose\_predictions.txt} per sequence with 8 columns
(\texttt{k tx ty tz qx qy qz qw}), one row per input frame, frame 0 the identity.

\subsection{iMED Novel View Synthesis (NVS)}

\paragraph{Hardware.} Development on a single NVIDIA Quadro RTX 8000 (48\,GB). The container requires
CUDA; a CPU-only build of the same method was rejected twice by the evaluator regardless of its wall-clock
time, so the entire warp, $z$-buffer, hole filling and photometric transform run on the GPU.

\paragraph{Training.} \textbf{No per-sequence optimisation and no training on iMED data.} The only learned
component is the instrument segmenter (disclosed below), which is used inference-only.

\paragraph{Pretrained models and external data (disclosure).} The instrument segmenter is a U-Net with a
ConvNeXt-Base encoder (\texttt{tu-convnext\_base.dinov3\_lvd1689m} via \texttt{timm}; a ConvNeXt-Base
distilled from DINOv3 features~\cite{liu2022convnext,simeoni2025dinov3}) built with
\texttt{segmentation\_models\_pytorch}~\cite{iakubovskii2019smp}. It was trained by our group for binary
instrument segmentation on the \textbf{public Cholec80 dataset}~\cite{twinanda2017endonet}; the checkpoint
predates this challenge, contains no iMED data, and is used here frozen. The weights are baked into the
image so the container runs with \texttt{-{}-network=none}. The photometric constants are the only other
fitted quantity, and they are fitted on the released public iMED NVS sequences (no hidden data, no
target-view ground truth at inference).

\paragraph{Segmenter inference settings.} BGR$\rightarrow$RGB, bilinear resize to $512\times512$, ImageNet
normalisation (mean $[0.485,0.456,0.406]$, std $[0.229,0.224,0.225]$), sigmoid output, bilinear resize of
the probability map back to the input resolution, threshold $0.30$. The threshold was selected on the
development set ($0.30$: 22.222; $0.50$ and $0.70$: 22.18), i.e.\ leaning slightly toward over-calling: a
missed instrument re-projects metal onto tissue, which costs more than inpainting a little extra tissue.

\paragraph{Warp settings.}
\begin{itemize}\itemsep1pt
  \item Full native RGB resolution ($1280\times1024$); \texttt{K.txt} is used verbatim, as it is
        calibrated at full resolution. Output resolution equals the target resolution exactly --- the
        official metric does not resize.
  \item Depth is read from \texttt{endoscope2/depthL/*.npy} and upsampled to RGB resolution with
        \textsc{nearest} interpolation, deliberately, so that no interpolated sample straddles a depth
        discontinuity.
  \item $z$-buffer: depth is quantised with a scale factor of 64 and packed as
        \texttt{(z << 32) | source\_index} into an \texttt{int64}; visibility is a single
        \texttt{scatter\_reduce\_(..., reduce="amin")} over the target pixel index. Nearest-integer
        (not bilinear) target coordinates.
  \item Hole filling: push--pull pyramid with 6 levels, $2\times$ average pooling with
        \texttt{ceil\_mode}, bilinear upsampling, and a validity threshold of $10^{-5}$ on the
        accumulated weight.
  \item Photometric transform, in BGR and $[0,1]$: $\hat{y}=\mathrm{clip}(g\odot x + b)$ with
        $g = [0.665668,\ 0.614038,\ 0.642004]$ and $b = [0.097009,\ 0.143732,\ 0.170179]$, fitted by
        least squares over the officially scored pixels of all 20 public sequences.
\end{itemize}

\paragraph{Graceful degradation.} If the segmenter cannot be constructed, the pipeline falls back to the
supplied \texttt{endoscope2/toolL} masks; if those are absent too, the frame is splatted with the
instrument included (the pre-exclusion behaviour). A frame with no depth map is written as black rather
than aborting the run.

\paragraph{Runtime.} $\approx$59\,s per sequence end to end (all 20 development sequences,
3{,}953 frames, on the RTX 8000); the warp itself is $\approx$79\,ms per frame. For reference, the CPU
implementation of the same algorithm cost 264\,s per sequence.

\paragraph{Container.} \texttt{nvidia/cuda:11.8.0-cudnn8-runtime-ubuntu22.04}; Python 3.10;
\texttt{torch 2.1.2+cu118}, \texttt{torchvision 0.16.2+cu118}, \texttt{numpy 1.26.4},
\texttt{scipy 1.11.4}, \texttt{opencv-python-headless 4.9.0.80}, \texttt{timm 1.0.24},
\texttt{segmentation-models-pytorch 0.5.0} (installed with \texttt{-{}-no-deps} so that its unpinned
\texttt{torch} requirement cannot silently replace the CUDA-11.8 build). Image size 10.1\,GB. Contract:
\texttt{/input} mounted read-only; output written as
\texttt{/output/<seq>/renders/00000.png}\,\ldots, one $1280\times1024$ 8-bit RGB PNG per source frame, in
sorted order. Ground-truth target views are never read.

% =============================================================================
\section{Results on Public Set (Optional)}
% =============================================================================

% This section is optional.
% Report only results measured on a released development/validation split or
% results explicitly supplied by the organizers. Do not infer hidden-test results.
%
% Clearly identify the split used.

\subsection{iMED Pose Estimation (PE)}

\textbf{Split: the 19 released test sequences that ship with ground-truth \texttt{pose.txt}}, scored with
the organiser evaluation script (Sim(3)-aligned ATE).

\begin{center}
\begin{tabular}{lccc}
\toprule
Method & mean ATE [mm] & Registered frames & Runtime \\
\midrule
Official baseline (ALIKED+LightGlue, essential) & 2.540 & 100\% & $\sim$14\,s/seq \\
\textbf{Tri3D-GridBA (submitted)}               & \textbf{0.953} & 100\% & $\sim$0.43\,s/frame \\
\bottomrule
\end{tabular}
\end{center}

Under the RMSE variant of the same metric, the submitted configuration scores $1.079$\,mm over the same
19 sequences and $2.070$\,mm over the 6 hardest of them. Runtimes are on a Quadro RTX 8000.

\textbf{Organiser-supplied validation leaderboard (6 sequences, \texttt{mean\_ate\_rmse\_mm}).} The
submitted image scored $2.296$\,mm. We report this only for completeness: a bootstrap over paired
per-sequence differences indicates a resolution of $\pm 0.10$\,mm on a 6-sequence set, so it should not be
compared at finer granularity than that. The gap between the leaderboard value and the development value
is consistent with the leaderboard subset being weighted toward the fast-circular sequences that we
identify above as an observability limit.

\textbf{Qualitative observation.} Errors are strongly session-structured rather than frame-structured.
On smooth trajectories the recovered translation tracks the ground truth closely and the grid BA barely
moves it; on fast circular motion the per-frame translation is systematically about twice too large even
when the rigid fit is well-conditioned by every internal diagnostic we could construct
(match count, scale variance, triangulation residual, disparity). Four methodologically independent
estimators fail on the same sequences, which is why we characterise this as a property of the data.

\subsection{iMED Novel View Synthesis (NVS)}

\textbf{Split: the 20 released \texttt{iMED\_NVS} sequences (3{,}953 frames) with endoscope-1 ground
truth}, scored with the official \texttt{metrics.py} inside the official
overlap\,$\times$\,instrument-free region.

\begin{center}
\begin{tabular}{lccc}
\toprule
Configuration (all 20 sequences, all frames, measured from the containers) & PSNR & SSIM (gray) & Runtime \\
\midrule
Warp $+$ push--pull $+$ poly2, instrument re-projected     & 22.3625 & 0.6892 & 31.5\,s/seq \\
\quad $+$ instrument exclusion                              & 22.4880 & 0.6939 & 44.3\,s/seq \\
\textbf{$+$ affine photometric transform (submitted)}       & \textbf{22.3897} & 0.6830 & 59.2\,s/seq \\
\bottomrule
\end{tabular}
\end{center}

For reference, on the single sequence for which we ran the official Endo-4DGS baseline to convergence
(\texttt{session\_004\_scene\_2\_tool\_1}, native resolution), the baseline scores PSNR 21.14 / SSIM 0.699
/ LPIPS 0.0190 after $\sim$2\,h of per-sequence optimisation, and our warp scores PSNR 21.45 / SSIM 0.721
/ LPIPS 0.0198 with no optimisation at all.

\textbf{Organiser-supplied validation leaderboard.} The submitted image scored PSNR $21.156$,
SSIM $0.621$, LPIPS $0.200$.

\textbf{Qualitative observations.} (i) Reconstruction quality inside the scored region is limited by
per-scene calibration consistency, not by the renderer: a sequence-constant image shift of 8--35\,px is
present in every sequence and accounts for $\sim$2\,dB, but its magnitude and direction are scene-specific
and do not transfer to unseen sessions. (ii) Disocclusion holes cover $19.6\%$ of the scored pixels and
are the second-largest term ($+0.99$\,dB if filled with ground truth); the visible failure mode of the
push--pull filler is a smear across occlusion boundaries, and the structure-aware alternative that fixes
this visually is the one that costs SSIM. (iii) The photometric transform is not purely a colour
correction: part of its gain is the MSE-optimal contraction toward the conditional mean, which raises PSNR
and SSIM while lowering contrast and worsening LPIPS. We accepted this because PSNR is the primary metric,
and we note it here so the effect is not mistaken for a genuine appearance model.

% Optional qualitative figure:
%
% \begin{figure}[h]
%     \centering
%     \includegraphics[width=0.9\linewidth]{qualitative.png}
%     \caption{\TODO{Brief description and data split used.}}
% \end{figure}

% =============================================================================
% cite as normal like 3D Gaussian Splatting~\cite{kerbl2023}

\bibliographystyle{unsrt}
\bibliography{references}

\end{document}
